% Characters the font has no glyph for, typed raw as in real papers: TeX reports each one missing and sets nothing for it (an OpenType font sets its .notdef), and the page draws nothing there either.
\documentclass{article}
\usepackage[T1]{fontenc}
\usepackage{lmodern}
% one page with nothing but the text: checked against plain lualatex too (cases.ts)
\pagestyle{empty}
\begin{document}
A paragraph in a T1 font, with characters typed raw that its 256 slots do
not hold as Unicode: the reference to Paweł Borówka of Łódź, set in T1, loses
its ł and Ł and ź, and a Chernoff bound typed with the ligature character loses
the ligature, while the letters around them stay where TeX put them. The same
words with the encoding's own commands, Pawe\l{} and \L\'od\'z, keep theirs.

{\fontencoding{TU}\selectfont
A paragraph in Latin Modern's OpenType fonts, which have no box-drawing
characters: a frame drawn in a listing as ╔══╗ and ║ and ╚══╝ is left out,
each character set as the font's empty .notdef glyph, which keeps the line's
spacing as TeX gave it. The text after it, in the same font, goes on as
usual, long enough to break into several lines at every width.\par}

A last paragraph in the T1 font again, with nothing missing, to show that
the lines after the two above are set where TeX set them.
\end{document}
